% Translation quotients, reduced DT series, and automorphic denominators.
% Primary anchors: Oberdieck--Pandharipande, Gritsenko--Nikulin,
% Borcherds, Lorgat 2020.

\providecommand{\kcat}{\kappa_{\mathrm{cat}}}
\providecommand{\kch}{\kappa_{\mathrm{ch}}}
\providecommand{\kBKM}{\kappa_{\mathrm{BKM}}}
\providecommand{\kfiber}{\kappa_{\mathrm{fiber}}}
\providecommand{\Zpartition}{Z}
\providecommand{\gBKM}{\mathfrak{g}}
\providecommand{\DeltaFive}{\Delta_5}
\providecommand{\PhiTen}{\Phi_{10}^{\mathrm{un}}}

\section*{Translation Quotients and Reduced DT Denominators}

\paragraph{Free translation quotient.}
Let \(S\) be a K3 surface with a symplectic automorphism \(g_N\) of
order \(N\leq 8\). Let \(E\) be an elliptic curve and let \(e_0\in E\)
be a torsion point of exact order \(N\). The diagonal action
\[
  \mathbb Z/N\mathbb Z\curvearrowright S\times E,
  \qquad
  (s,e)\longmapsto(g_Ns,e+e_0)
\]
is free because translation by \(e_0\) has no fixed points. It
preserves the holomorphic \(3\)-form, hence
\[
  X_N=(S\times E)/(\mathbb Z/N\mathbb Z)
\]
is a smooth Calabi--Yau threefold whenever \(g_N\) exists.

This translation quotient is not the same object as the stricter
origin-fixing elliptic-automorphism lane used for the primitive
\(K3\times E\) CHL denominator table. The latter has the five levels
\[
  N\in\{1,2,3,4,6\},
  \qquad
  c_N(0)=(10,8,6,4,2),
  \qquad
  \kBKM(\Phi_N)=c_N(0)/2=(5,4,3,2,1).
\]
Orders \(5,7,8\) can appear in the translation-quotient geometry, but
they are not additional rows of this five-level denominator table.

\paragraph{Curve classes.}
Assume \(S\to\mathbb P^1\) is elliptically fibered with zero section
\(s_1\) and an \(N\)-torsion section \(s_2\). The translates
\[
  s_k=s_{k-1}+s_2,\qquad 1\leq k\leq N,
\]
define \(N\) sections. With \(F\in\operatorname{Pic}(S)\) the fibre
class, set
\[
  \beta_h=\frac1N(s_1+s_2+\cdots+s_N+hF),
  \qquad h\geq0.
\]
This class belongs to the quotient lattice after descent to \(X_N\).
The divisibility by \(N\) is a descent statement, not a statement about
individual section classes on \(S\).

\paragraph{Reduced DT series.}
For curve class \((\beta_h,d)\), define
\[
  \Zpartition^{X_N}(q,t,p)
  =
  \sum_{h\geq0}\sum_{d\geq0}\sum_{n\in\mathbb Z}
  \mathrm{DT}_{n,(\beta_h,d)}\,
  q^{d-1}t^{\langle\beta_h,\beta_h\rangle/2}(-p)^n .
\]
The reduced invariant is the Behrend-weighted Euler characteristic of
the Hilbert scheme modulo the elliptic translation action:
\[
  \mathrm{DT}_{n,(\beta_h,d)}
  =
  \int_{\operatorname{Hilb}^{n}(X_N,\beta_h)/E}\nu_{\mathrm{Behr}} .
\]
The monomial \(q^{d-1}t^{\langle\beta_h,\beta_h\rangle/2}(-p)^n\)
fixes the DT convention. Changing this convention changes the displayed
automorphic scalar by an explicit monomial, not by a theorem.

\paragraph{The \(N=1\) scalar.}
For \(X_1=S\times E\), the reduced DT partition function is the
Oberdieck--Pandharipande scalar
\[
  \Zpartition^{S\times E}
  =
  -D_5^{-2}
  =
  -4096\,\DeltaFive^{-2},
  \qquad
  D_5=64^{-1}\DeltaFive .
\]
Equivalently, in the unnormalised Igusa convention,
\[
  \PhiTen=\DeltaFive^2 .
\]
The K3 elliptic genus is
\[
  \operatorname{Ell}(K3;\tau,z)=2\phi_{0,1}^{\mathrm{EZ}}(\tau,z).
\]
The primitive Borcherds denominator uses
\(\phi_{0,1}^{\mathrm{EZ}}\) and gives
\[
  \kBKM(\DeltaFive)=c_1(0)/2=10/2=5.
\]
The doubled K3 elliptic genus gives the square
\(\DeltaFive^2\), of weight \(10\). Thus the reduced-DT scalar is
square-level evidence. It does not by itself construct the primitive
BKM root datum or a compact Hall product.

\paragraph{Twisted denominators.}
A twisted reduced-DT series
\[
  \Zpartition^{X_N}_{L,h_M}(q,t,p)
\]
may produce an automorphic denominator after one supplies:
\[
  \text{the lattice polarisation }L,\quad
  \text{the commuting pair }(g_N,h_M),\quad
  \text{the Jacobi input},\quad
  \text{the product normalisation},\quad
  \text{the CY host and orientation data}.
\]
Only then can one assert an identity of the form
\[
  -\bigl(\Zpartition^{X_N}_{L,h_M}\bigr)^{-1}
  =
  \text{Borcherds denominator}.
\]
The associated Lie superalgebra
\(\gBKM_{(g_N,h_M)}\) requires a root lattice, Weyl chamber, Weyl
vector, parity split, and root-multiplicity rule. A scalar modular form
is not a Lie-superalgebra construction.

\paragraph{Jacobi basis.}
The standard weak Jacobi basis used in the twining formulas is
\[
  A=\frac14\phi_{0,1},\qquad
  B=\phi_{-2,1}.
\]
The symbol \(A\) is not \(\frac14\phi_{0,1}^2\). The square belongs to
the Igusa/DVV scalar lane, not to the weight-zero index-one Jacobi
basis.

\paragraph{Scope.}
The translation quotient gives a broad family of smooth CY\(_3\)'s.
The five-level primitive CHL denominator table is a narrower
automorphic lane. The \(N=1\) reduced-DT scalar matches
\(\Delta_5^{-2}\) after the stated normalization; the primitive BKM
weight remains
\[
  \kBKM(\DeltaFive)=5.
\]
For \(N>1\), the automorphic denominator, DT scalar, and BKM algebra
are three distinct levels of structure. Identifying them requires
comparison maps, not only matching constants.
